\subsection{Policy dimension}
The latent factor labels are not identified.  If the raw characteristic map is
$\phi(z)\in\mathbb R^K$, any nonsingular change of coordinates can be offset by
an inverse change in the state coefficients without changing stock weights.
We therefore define dimension from the portfolio rule itself rather than from
neural units.

Let $D_t$ be the diagonal matrix of cross-sectional root-mean-square scores in
\eqref{eq:score_normalization}.  Because $\widetilde\phi_t=D_t^{-1}\phi$, the
stock-weight rule can be written
\begin{equation}
N_t w_{i,t}=\phi(z_{i,t})'c_t,\qquad c_t=D_t^{-1}b_t.
\label{eq:separable_policy}
\end{equation}
Thus the normalization used for numerical identification does not alter the
separable structure: it merely moves a time-varying scale from the
characteristic map into the state coefficient.
Define the empirical characteristic and state Gram matrices
\begin{equation}
G_\phi=\mathbb E_t\!\left[\frac{1}{N_t}\sum_{i=1}^{N_t}
\phi(z_{i,t})\phi(z_{i,t})'\right],
\qquad
G_c=\mathbb E_t[c_tc_t'].
\label{eq:policy_grams}
\end{equation}
The nonzero eigenvalues of
\begin{equation}
\mathcal P=G_\phi^{1/2}G_cG_\phi^{1/2}
\label{eq:policy_operator}
\end{equation}
are the squared singular values of the learned separable portfolio rule under
the empirical product reference measure.  Write them as
$\lambda_1\geq\cdots\geq\lambda_K\geq0$ and call them the \emph{policy
spectrum}.  The spectrum measures how many characteristic-by-state directions
are needed to represent the learned stock-weight function; it does not count
equilibrium risk factors in the market.
\begin{proposition}[Invariance of the policy spectrum]
The eigenvalues $\{\lambda_j\}_{j=1}^K$ are unchanged by any nonsingular
reparameterization $\widetilde\phi=A\phi$ and
$\widetilde c=A^{-\prime}c$.
\end{proposition}
\begin{proof}
The transformed Gram matrices are
$\widetilde G_\phi=AG_\phi A'$ and
$\widetilde G_c=A^{-\prime}G_cA^{-1}$.
The eigenvalues of \eqref{eq:policy_operator} equal the nonzero eigenvalues of
$G_\phi G_c$.  Under the transformation,
\[
\widetilde G_\phi\widetilde G_c
=A(G_\phi G_c)A^{-1},
\]
which is similar to $G_\phi G_c$ and therefore has the same spectrum.
\end{proof}

Normalize $q_j=\lambda_j/\sum_{\ell}\lambda_\ell$.  Our primary measure of
policy dimension is the entropy rank
\begin{equation}
K_{\rm policy}=\exp\!\left(-\sum_jq_j\log q_j\right),
\label{eq:policy_rank}
\end{equation}
with the participation ratio reported as a robustness measure.
This object also has a direct approximation interpretation.  By the Schmidt
expansion, the squared $L^2$ error of the best rank-$r$ separable
approximation to the learned policy is proportional to
$\sum_{j>r}\lambda_j$.  Hence a concentrated policy spectrum means that a
rich neural network uses its flexibility to discover nonlinear functions
inside a low-rank characteristic-state policy, not that the underlying return
panel itself must have low rank.

\subsection{Canonical traded factors}
Policy rank and traded-factor interpretation are related but distinct.  For
economic interpretation we work with the returns of the $K$ learned managed
portfolios in \eqref{eq:learned_factor}.  Let
\begin{equation}
S=\mathbb E[F_{t+1}F_{t+1}'],\qquad
g_{t+1}=S^{-1/2}F_{t+1},\qquad d_t=S^{1/2}b_t.
\label{eq:risk_standardization}
\end{equation}
Then $R^p_{t+1}=d_t'g_{t+1}$ and $\mathbb E[g_{t+1}g_{t+1}']=I$.
We diagonalize
\begin{equation}
H=\mathbb E[d_td_t']
\label{eq:economic_operator}
\end{equation}
and rotate $(g_{t+1},d_t)$ by the corresponding orthogonal eigenvectors.
The resulting canonical traded factors are ordered by risk-standardized
allocation intensity rather than by arbitrary hidden-unit labels.
The eigenvalues of $H$ are invariant to nonsingular rotations of the latent
factor coordinates by the same whitening argument.  We refer to their
entropy rank as the \emph{allocation rank}.  It measures concentration of the
estimated strategy across risk-standardized traded directions; unlike
$K_{\rm policy}$, it is not our definition of functional policy dimension.
In particular, a static allocator mechanically collapses the implemented
portfolio to one characteristic score even if the latent factor return space
has many dimensions.

The economically decisive diagnostic is therefore out-of-sample truncation.
Denote the $j$th canonical factor by $G_{j,t+1}$ and its allocation by
$D_{j,t}$. Define
\begin{equation}
R^{p,(r)}_{t+1}=\sum_{j=1}^{r}D_{j,t}G_{j,t+1}.
\label{eq:truncated_portfolio}
\end{equation}
We report the Sharpe ratio of $R^{p,(r)}$, its correlation with the full
portfolio, and the relative payoff mean-squared error as $r$ increases.  If a
small number of directions reproduces the full strategy out of sample, low
economic dimension is a statement about portfolio value rather than a visual
feature of an eigenvalue plot.
